\chapter{Allodynia}

\section*{Definition}
Allodynia is pain caused by a stimulus that does not normally provoke pain, such as light touch, gentle pressure, or mild warmth or cold. It may be mechanical or thermal and can occur with neuropathic pain, migraine, or sensitized tissues. It is distinct from hyperalgesia, which is an increased response to a stimulus that is normally painful.

\section*{When to See a Doctor}

\begin{itemize}
 \item Pain provoked by light touch, clothing, or temperature changes
 \item New, spreading, progressive, or disabling pain, especially with numbness or weakness
 \item Associated with skin color changes, edema, or sweating (suggesting CRPS)
 \item Progressive or disabling pain not responsive to over-the-counter analgesics
\end{itemize}

\section*{How It Develops}
Allodynia can result from increased sensitivity in nerves and pain-processing pathways after injury, inflammation, infection, or repeated pain signals. Peripheral sensitization lowers the activation threshold of nerve endings; central sensitization increases the response of spinal and brain networks. The exact mechanisms vary by condition and are still being studied, so a symptom does not identify one specific disease.

\section*{Causes}
\begin{itemize}
 \item \textbf{Neuropathic:} Diabetic peripheral neuropathy, postherpetic neuralgia, trigeminal neuralgia, Fabry disease
 \item \textbf{Central sensitization syndromes:} Fibromyalgia, complex regional pain syndrome (CRPS), migraine
 \item \textbf{Spinal cord injury:} Central pain syndrome, syringomyelia
 \item \textbf{Post-surgical:} Post-mastectomy pain, post-thoracotomy pain, phantom limb pain
 \item \textbf{Infectious:} HIV-associated neuropathy, postherpetic neuralgia induced by varicella-zoster virus
 \item \textbf{Demyelinating:} Multiple sclerosis, Guillain-Barr\'{e} syndrome
 \item \textbf{Chemotherapy-induced:} Platinum-based drugs, taxanes, bortezomib, thalidomide
 \item \textbf{Toxins:} Alcohol-related neuropathy, heavy metal toxicity
\end{itemize}

\section*{Related Symptoms}
\begin{itemize}
 \item Pain with light brushing, clothing contact, or wind (mechanical allodynia)
 \item Pain with cold or warm stimuli not typically painful (thermal allodynia)
 \item Hyperalgesia (increased pain to normally painful stimuli)
 \item Paresthesias, dysesthesias, or burning sensation
 \item Autonomic changes: skin color, temperature, sweating, edema
\end{itemize}
\textbf{Important:} Allodynia is a sensory finding, not a diagnosis. Its presence does not by itself prove central sensitization or predict which medicine will work.

\section*{Medical Evaluation}
\begin{itemize}
 \item History and examination to map the pain, identify sensory loss or weakness, assess skin temperature and color, and look for signs of nerve, spinal cord, joint, or vascular disease
 \item Quantitative sensory testing or skin biopsy only when a specialist needs to evaluate suspected small-fiber neuropathy; these tests are not required for every case
 \item Nerve conduction studies and electromyography when large-fiber neuropathy, radiculopathy, or another nerve disorder is suspected
 \item MRI of the brain or spine when examination and symptoms suggest a central lesion, such as spinal cord compression or multiple sclerosis
 \item Targeted laboratory tests, often including glucose or HbA1c, vitamin B12, kidney function, and thyroid testing; additional tests for inflammation, infection, or autoimmune disease depend on the history and examination
\end{itemize}

\section*{Treatment Options}
\begin{itemize}
 \item Treat the underlying cause when identified. For neuropathic pain, clinicians may offer amitriptyline, duloxetine, gabapentin, or pregabalin and review benefit and adverse effects regularly.
 \item Topical lidocaine or capsaicin may help localized peripheral neuropathic pain. High-concentration capsaicin patches require professional application; compounded creams have uncertain evidence and quality.
 \item Trigeminal neuralgia requires condition-specific treatment such as carbamazepine or oxcarbazepine. Tramadol is generally reserved for short-term rescue when other options are unsuitable; routine long-term opioid treatment is usually avoided.
 \item Interventional treatments such as sympathetic blocks, spinal cord stimulation, or dorsal root ganglion stimulation may be considered by specialist teams for selected refractory conditions, not for uncomplicated allodynia
 \item Physical and occupational therapy with graded motor imagery and desensitization
 \item Cognitive-behavioral therapy and pain psychology for central sensitization syndrome
\end{itemize}

\section*{Self-Care and Home Management}
\begin{itemize}
 \item Protective clothing to minimize skin irritation from light touch
 \item Avoidance of known triggers (e.g., tight clothing, temperature extremes)
 \item Desensitization therapy with graded exposure to textures and temperatures
 \item Ask a pharmacist or clinician before using topical lidocaine or capsaicin, particularly on broken skin or with other local anesthetics; stop if severe irritation develops
 \item Regular sleep hygiene and stress reduction to lower central excitability
\end{itemize}

\section*{References}
\begin{thebibliography}{99}
\bibitem{allodynia:ref1} Jensen TS, Finnerup NB. ``Allodynia and hyperalgesia in neuropathic pain.'' \textit{Lancet Neurol}. 2014;13(9):924-935.
\bibitem{allodynia:ref2} Finnerup NB, Attal N, et al. ``Pharmacotherapy for neuropathic pain in adults: a systematic review.'' \textit{Lancet Neurol}. 2015;14(2):162-173.
\bibitem{allodynia:ref3} Woolf CJ. ``Central sensitization: implications for the diagnosis and treatment of pain.'' \textit{Pain}. 2011;152(3 Suppl):S2-S15.
\bibitem{allodynia:ref4} Ballantyne JC, Cousins MJ, et al. \textit{Bonica\'s Management of Pain}, 5th ed. Wolters Kluwer, 2018.
\bibitem{allodynia:ref5} UpToDate. ``Evaluation and management of allodynia.'' Wolters Kluwer, 2023.
\bibitem{allodynia:ref6} International Association for the Study of Pain. Allodynia and hyperalgesia in neuropathic pain. Updated 2020.
\bibitem{allodynia:ref7} National Institute for Health and Care Excellence. Neuropathic pain in adults: pharmacological management in non-specialist settings (CG173). Updated 2020.
\end{thebibliography}
